% Background and motivation (Week 23). Sources: reading/levy2020narya.md, reading/tirmazi2020borg.md,
% reading/resource_central2017.md, docs/v2/v1_final_state.md, docs/v2/v1_anomaly.md.
\section{Background and motivation}
\label{sec:background}

\paragraph{Predictive mitigation in production.}
Narya~\cite{levy2020narya} is the most detailed public account of predictive failure mitigation
for VMs. A rule set and a gradient-boosted model flag nodes likely to fail within days; for each
flagged node an online bandit chooses among composite actions such as marking the node
unallocatable, live-migrating its VMs, or sending it to investigation, and learns from the
interruptions observed over the following weeks. Its decision is made per node and prediction
rule. The VMs on a node are not ranked against VMs on other nodes, and what a migration costs in
network or memory enters only through rate limits, not through the choice. Narya reports its gains
relative to a fixed action plan, not to a tuned threshold, and only as percentages of
proprietary totals. Disk-error predictors at the same operator rank disks and act on the top few
under ``constraints of cost and capacity''~\cite{xu2018cdef}, again without saying which VMs to
move.

\paragraph{What public traces say about the inputs.}
Three properties of public data shape the question. First, failure signals are weak at the top of
the ranking: on the Backblaze fleet our calibrated 30-day model flags drives with a precision of
about 5\% in its top 1\% (Section~\ref{sec:model}), so most migrations protect VMs whose host
would not have failed. Second, VMs are short-lived and unequal: most Azure VMs live minutes, while
the long-lived minority holds almost all memory-time~\cite{cortez2017resource}, so the cost of an
interruption and the memory a migration moves vary by orders of magnitude across VMs. Third,
cluster schedulers place work for packing, priority and spread~\cite{tirmazi2020borg,hadary2020protean},
not for hardware risk, so there is no reason to expect critical work to sit on safer or riskier
machines, but also no published measurement.

\paragraph{Why the correlation was thought to matter.}
\label{sec:v1}
Our course prototype simulated \VVOneHosts{} Backblaze drives as hosts carrying Azure VMs, and
compared a cost-aware rule (migrate when $p\,c$ exceeds the migration cost) with a probability
threshold matched to the same number of migrations. With VMs placed at random it found the
cost-aware rule \VVOneAdvRandom\% behind; as placement was skewed so that critical VMs sat on
riskier hosts, the sign flipped, reaching \VVOneAdvTop\% at a rank correlation of \VVOneRhoTop. Read at face value
this says that workload awareness pays only above a break-even correlation, which makes the
correlation in real placement the quantity that decides the method's value.

A pre-registered diagnosis of the prototype~(eight hypotheses, each with a rule written before it
was tested) found otherwise. The threshold had received a larger budget: matched on migration
count, it moved \VVOneExtraHosts\% more hosts and \VVOneExtraGB\% more memory. And every migration was credited with any
failure of its host over the whole run, a median of \VVOneMedianLeadDays{} days after a decision
made on a 30-day probability, so the comparison rewarded predicting the distant future, which
neither rule was built to do. With the budget matched in memory and the accounting window equal
to the decision horizon, the cost-aware rule is \VVOneAdvAligned\% ahead at zero correlation.
The small size of that gain is itself predicted by the theory of Section~\ref{sec:approach} once
it is diluted by the share of total cost a 30-day decision can affect.

The prototype therefore left two questions open, which this paper answers in order: where real
placement sits on the correlation axis, and how large the advantage of budget-aware ranking is
there, with the accounting done right and the comparator tuned to use the same budget.
